\section{Background}
\label{sec:background}
In this paper,
we will use the following notation.
An overline indicates
logical states and operator representatives
e.g.\ $\ket{\overline{0}}, \overline{Z}$.
Physical operators are without overline
and the standard ones are typed upright:
I, X, Y, Z,
$\S =\diag(1, \i)$,
$\T =\diag(1, \e^{\i \uppi /4})$,
CX.
We use the shorthand $\H_\pm :=\frac{1}{\sqrt{2}}(\X \pm \Y)$.
Sets are typed calligraphically e.g.\ $\mathcal{Q}$.
If clear from the context,
we refer to elements from the Pauli (Clifford) group
simply as Paulis (Cliffords).
Given a 1-qubit unitary $U$,
define $U_\mathcal{Q} :=\prod_{q \in \mathcal{Q}} U_q$
for integer set $\mathcal{Q} \subseteq `{1, \dots, n}$,
where $U_q$ denotes $U$ acting on qubit $q$
and $\I$ acting on the other qubits.
The \emph{support} of a multiqubit unitary
is the set of qubits on which it acts something other than $\I$,
e.g., the support of $U_\mathcal{Q}$ is $\mathcal{Q}$.
A \emph{stabiliser circuit} is a quantum circuit comprising
Clifford gates and 1-qubit Z preparations and measurements~\cite{Aaronson2004}.
The parameters of a quantum error correction code are written as
$\llbracket n, k, d \rrbracket$
for physical qubit count $n$,
logical qubit count $k$,
and distance $d$.
In the case distance is unimportant,
we write $\llbracket n, k\rrbracket$.
The logical action of any operator $U$
on a code may be realised by
multiple physically distinct representatives;
throughout, we use $\overline{U}$ to denote
a \emph{specific} representative.

The set
$\mathcal P_k:=\{\I,\X,\Y,\Z\}^{\otimes k}$
of phaseless $k$-qubit Paulis
is a basis for the space of
$2^k \times 2^k$ Hermitian matrices over $\mathbb R$.
So,
the density operator of any $k$-qubit state $\rho$
can be decomposed into this basis:
\begin{equation}
\rho =
2^{-k} \sum_{L \in \mathcal{P}_k} \alpha_L L,
\qquad \alpha_L \in \mathbb R,
\qquad \alpha_\I =1.
\label{eq:pauli_decomposition}
\end{equation}
Here $\alpha_\I$ is obtained
by ensuring $\tr \rho =1$,
and noting $\tr L =2^k$ if $L =\I$ else 0.

Cliffords map Paulis to Paulis.
Any Clifford can be defined up to a global phase
by how it maps the X and Z operators on each input, e.g.:
\begin{equation}\label{eq:h_pm_action}
	\H_\pm: \begin{cases}
		\X &\mapsto \pm \i \X\Z \\
		\Z &\mapsto -\Z.
	\end{cases}
\end{equation}

\subsection{The Hexagonal Colour Code}
Here we review the $\llbracket n, 1, d \rrbracket$ hexagonal colour code
for its relevance to MSC,
and for examples later on in this paper.
For this code,
choose the representatives
$\overline{X} =\X^{\otimes n}$ and $\overline{Z} =\Z^{\otimes n}$.
For the $\llbracket 7, 1, 3 \rrbracket$ instance in particular,
we label the data qubits according to \cref{fig:d3_layout}
and choose the following set $\mathcal G$ of six stabiliser generators:
\begin{equation}
(G_1,\ldots,G_6) =(
\mathrm{XXXIXII},
\mathrm{IXXXIXI},
\mathrm{IIXIXXX},
\mathrm{ZZZIZII},
\mathrm{IZZZIZI},
\mathrm{IIZIZZZ}).
\label{eq:d3_color_stabilizer_generators}
\end{equation}
\begin{figure}[b]
	\begin{subfigure}{0.48\linewidth}
		\centering
		\includegraphics[width=0.55\textwidth]{figures/d3_layout.png}
		\caption{The physical qubit layout used in
		Stage 1 of the magic state cultivation protocol.
		Numbers 0--6 label the data qubits of
		the $\llbracket 7, 1, 3 \rrbracket$ colour code.
		Letters a--f label the ancilla qubits used for
		stabiliser and logical Clifford measurement.}
		\label{fig:d3_layout}
	\end{subfigure}
	\hfill
	\begin{subfigure}{0.48\linewidth}
		\centering
		\includegraphics[width=0.7\textwidth]{figures/BipartiteGraph.pdf}
		\caption{Representing the
		$\llbracket 19, 1, 5 \rrbracket$ colour code
		as a bipartite graph
		between light and dark nodes.
		Figure from \cite{Kubica2015}.}
		\label{fig:bipartite_graph}
	\end{subfigure}
	\caption{Distance-3 and -5 hexagonal colour codes.
	Each colour-shaded polygon
	represents an X- and a Z-type stabiliser generator,
	whose support is the set of qubits at its corners.}
\end{figure}
\noindent
We later use the following.
\begin{lem}\label{lem:logical_H_is_transversal}
On the $\llbracket n, 1, d \rrbracket$ hexagonal colour code,
logical $\H_\pm$ can be implemented transversally as follows.
Represent the physical qubits of the code as nodes,
and its plaquette edges as edges,
on a bipartite graph as the example in \cref{fig:bipartite_graph} shows.
Let $\mathcal{L}$ and $\mathcal{S}$ denote the two parts of the bipartition,
labelled such that $|\mathcal L|-|\mathcal S|\equiv 1 \pmod 4$;
such a labelling is always possible because $n$ is odd.
Then the implementation is
\begin{equation}
\overline{H_\pm} :=(\H_\pm)_\mathcal{L} (\H_\mp)_\mathcal{S}.
\label{eq:logical_H_is_transversal}
\end{equation}
\end{lem}
\begin{proof}
	It suffices to show $\overline{H_\pm}$ preserves the stabiliser group
	and has the desired logical action.
	The code comprises $w$-sided plaquettes where $w \in `{4, 6}$,
	each corresponding to a weight-$w$ stabiliser generator
	supported on an equal number of qubits from $\mathcal{L}$ and $\mathcal{S}$.
	Given this and \cref{eq:h_pm_action},
	$\overline{H_\pm}$ transforms the stabiliser generators as
	$\X^{\otimes w} \mapsto \X^{\otimes w} \Z^{\otimes w}$
	or $\Z^{\otimes w} \mapsto \Z^{\otimes w}$,
	hence preserves the stabiliser group.
	Also,
	$\overline{H_\pm}$ maps the logical X and Z observables in the desired way:
	$\X^{\otimes n}
	\mapsto\pm (-1)^{|\mathcal S|} \i^n \X^{\otimes n} \Z^{\otimes n}
	=\pm \i \X^{\otimes n} \Z^{\otimes n}$
	and $\Z^{\otimes n} \mapsto -\Z^{\otimes n}$.
\end{proof}
\noindent
This bipartition-based construction
was originally used for $\overline{S}$
\cite[\S II.C]{Kubica2015}
(hence provides a proof shortcut using $\overline{H_+} =\e^{-\i \uppi/4} \overline{SX}$),
but note other valid allocations of $\H_+, \H_-$ exist
\cite[Lemma~3]{Tansuwannont2025}.
For the rest of this paper,
$\overline{H_+}$ for hexagonal colour codes will refer to
the implementation in \cref{eq:logical_H_is_transversal}.

\subsection{Fault Distance}\label{sec:fault_distance}
In quantum error correction we naturally deal with
noisy quantum circuits that propagate and manipulate logical information.
These circuits usually contain measurements
that enable the construction of detectors:
a \emph{detector} is a set of Pauli measurements
whose parity is deterministic in the absence of noise.
For a given noisy circuit run,
the \emph{syndrome} is a binary vector $\v s$ whose support
indicates the detectors whose parities differ from expected.
We call $\v s =\v 0$ a \emph{trivial} syndrome.
We define an \emph{error event} as an unintended event
that occurs independently (from other error events)
at a specific location in the circuit
with a probability proportional to
some characteristic noise level $p$ of the system
e.g.\ a Pauli X after an S gate.
A \emph{hook error} is an error event that,
due to a physical entangling gate in the circuit,
spreads to eventually affect more than one data qubit.

Often, different error events
	flip the same set of detectors
	and have the same resultant Pauli effect on the data qubits
		when propagated to the end of the stabiliser circuit.
Whenever this occurs,
we group these equivalent error events together
into the same \emph{fault},
defined abstractly by that detector signature and resultant effect
(we give a proper definition in \cref{sec:faults}).
Then, instead of modelling noise as a large set of error events,
we can model it more efficiently as a smaller set of independent faults.
For those familiar with the \emph{detector error model} formalism~\cite{Gidney2021a}:
a fault is similar to an error mechanism
but carries slightly more information;
we define it this way on purpose to allow
subsequent propagation of the resultant effect through non-Cliffords.
A \emph{$k$-fault configuration} is a set of $k$ faults occurring together;
its \emph{fault count} is its cardinality, $k$.

\emph{Fault distance},
also known as \emph{circuit distance},
generalises the notion of code distance
(a property of the code) to circuits.
It is the minimum number of error events
(equivalently, faults)
needed to flip a logical observable
without flipping any detectors \cite[\S 3.1]{Gidney2024}.
For a subcircuit with quantum inputs,
we define its fault distance assuming error-free inputs
and count only error events occurring within the subcircuit.
A high fault distance is desirable because
as $p \to 0$, and conditioned on a trivial syndrome,
the logical error rate of a circuit of fault distance $d$ is $\upTheta(p^d)$.

\subsection{Magic State Cultivation}
\label{sec:magic_state_cultivation}

\begin{figure}
	\centering
	\begin{tikzpicture}
\begin{yquant}[/yquant/operator/minimum width=0mm]

qubit {} d[3];
nobit a[3];
nobit f;

init {$\ket{\overline{T}}$} (d);
init {$\ket{+}$} a[0];
init {$\ket{0}$} a[1,2];

align -;
cnot a[1]|a[0];
cnot a[2]|a[0];
init {$\ket{0}$} f;
cnot f|a[1];
align -;
cnot f|a[2];
dmeter {Z} f;
discard f;

box {$\H_\pm$} d[0]|a[0];
[operator/separation=-2.5mm]
box {$\H_\pm$} d[1]|a[1];
[operator/separation=-2.5mm]
box {$\H_\pm$} d[2]|a[2];

dmeter {X} a;
discard a;

barrier (-);

init {$\ket{+}$} a[0];
init {$\ket{0}$} a[1,2];

cnot a[1]|a[0];
cnot a[2]|a[0];

box {$\H_\pm$} d[0]|a[0];
[operator/separation=-2.5mm]
box {$\H_\pm$} d[1]|a[1];
[operator/separation=-2.5mm]
box {$\H_\pm$} d[2]|a[2];

cnot a[2]|a[0];
cnot a[1]|a[0];

dmeter {X} a[0];
dmeter {Z} a[1,2];
discard a;

barrier (-);

init {$\ket{+}$} a;

box {$\H_\pm$} d[0]|a[0];
[operator/separation=-2.5mm]
box {$\H_\pm$} d[1]|a[1];
[operator/separation=-2.5mm]
box {$\H_\pm$} d[2]|a[2];

cnot a[2]|a[0];
cnot a[1]|a[0];

dmeter {X} a[0];
discard a[0];
init {$\ket{+}$} a[0];

cnot a[1]|a[0];
cnot a[2]|a[0];

box {$\H_\pm$} d[2]|a[2];
[operator/separation=-2.5mm]
box {$\H_\pm$} d[1]|a[1];
[operator/separation=-2.5mm]
box {$\H_\pm$} d[0]|a[0];

dmeter {X} a;
discard a;

\end{yquant}
\end{tikzpicture}
	\begin{subfigure}{0.25\linewidth}
		\caption{The \emph{short single check}, used in
		Itogawa et al.~\cite{Itogawa2025}.
		This involves preparing a fault-tolerant GHZ state
		and measuring it transversally;
		the $\overline{H_+}$ measurement is the
		parity of these X measurements.}
		\label{fig:short_single_check_sketch}
	\end{subfigure}
	\hfill
	\begin{subfigure}{0.24\linewidth}
		\caption{The \emph{long single check}, used in
		Sahay et al.~\cite{Sahay2026}.
		This involves encoding a GHZ state,
		unencoding it,
		and measuring it;
		the $\overline{H_+}$ measurement is the
		single physical X measurement.}
		\label{fig:long_single_check_sketch}
	\end{subfigure}
	\hfill
	\begin{subfigure}{0.45\linewidth}
		\caption{The \emph{double check}, used in
		Gidney et al.~\cite{Gidney2024}.
		This is conceptually (a) preceded by its time reverse,
		and constitutes two $\overline{H_+}$ measurements:
		the first is the X measurement midway;
		the second is the parity of the X measurements at the end.}
		\label{fig:double_check_sketch}
	\end{subfigure}
	\caption{A sketch of three different methods to measure
	$\overline{H_+}$ from \cref{lem:logical_H_is_transversal}.
	The top three qubits represent the $n >3$ data qubits of the colour code.
	The size of the GHZ state in each method is flexible~\cite[\S B.3.b]{Sahay2026}.}
	\label{fig:check_options}
\end{figure}
There are two variants of MSC:
MSC-3 and MSC-5,
named after the distance of the largest colour code used.
Both variants comprise three stages:
injection, cultivation, and escape.

We first describe MSC-3.
The magic state $\ket{\T} \propto \ket{0} +\e^{\i \uppi/4} \ket{1}$
is the state stabilised by the Clifford $\H_+$,
so $\ket{\overline{T}}$ is stabilised by $\overline{H_+}$.
\begin{description}
	\item[Stage 1] (injection) prepares a $\ket{\overline{T}}$ state
	in the $\llbracket 7, 1, 3 \rrbracket$ colour code
	using a unitary circuit of fault distance 1,
	followed by one stabiliser measurement round.
	\item[Stage 2] (cultivation) measures the observable $\overline{H_+}$ twice;
	out of the possible methods outlined in \cref{fig:check_options} for this,
	MSC uses the \emph{double check} in \cref{fig:double_check_sketch}
	as it is the most amenable to 2D nearest-neighbour connectivity.
	Its explicit form for MSC-3,
	which we call the \emph{distance-3 double check},
	is shown in \cref{fig:d3a6};
	this relies on the decompositions
	$\H_+ =\T\X\T^\dag$ and
	$\H_- =\T^\dag \X\T$,
	and has two further optimisations:
	$\T^\dag$ cancels $\T$ between the two applications of controlled-$\H_\pm$,
	and the GHZ-state (un)encoding is morphed onto some data qubits
	for connectivity reasons.
	Note the circuit in the original paper~\cite[Figure~7]{Gidney2024} mistakenly
	measures $\H_+^{\otimes 7}$ which is logically equivalent to $\overline{H_-}$.
	\item[Stage 3] (escape) converts the colour code into a much larger surface code.
	Throughout the whole protocol up until this point:
	any detector flip causes the attempt to be aborted.
	Otherwise, after conversion,
	some number of stabiliser measurement rounds are performed;
	the syndrome is given to a decoder and the state is accepted
	if the confidence of the decoder exceeds some preset threshold.
\end{description}
MSC-3 is summarised in \cref{fig:stage_flowchart},
and is intended to have fault distance 3 before its escape stage.
\begin{SCfigure}
	\includegraphics{figures/from_yquant/d3a6.pdf}
	\caption{The double-check circuit of distance-3 T-state cultivation.
	The qubits are labelled according to \cref{fig:d3_layout}.
	Each X measurement corresponds to a detector.
	The dashed box is the region shown in \cref{fig:2_fault_configuration}.
	For S-state cultivation, replace T with S;
	click
	\href{https://algassert.com/crumble\#circuit=Q(0,0)0;Q(0,1)1;Q(1,0)2;Q(1,1)3;Q(2,0)4;Q(2,1)5;Q(2,2)6;Q(2,3)7;Q(3,0)8;Q(3,1)9;Q(3,2)10;Q(4,0)11;Q(4,1)12;POLYGON(0,0,1,0.25)8_11_10_5;POLYGON(0,1,0,0.25)7_10_5_3;POLYGON(1,0,0,0.25)3_5_8_0;TICK;MPP_Y0*Y8*Y11*Y10*Y7*Y5*Y3;TICK;MPP_X8*X0*X3*X5;TICK;MPP_X3*X5*X10*X7;TICK;MPP_X11*X8*X5*X10;TICK;MPP_Z8*Z0*Z3*Z5;TICK;MPP_Z7*Z10*Z5*Z3;TICK;MPP_Z11*Z8*Z5*Z10;TICK;RX_12_9_4_2_6_1;S_3_8_10;S_DAG_0_5_7_11;TICK;CX_1_0_9_8_6_7_4_5_2_3_12_11;TICK;CX_3_1_5_6_9_12;TICK;CX_5_3_9_10;TICK;CX_5_9;TICK;MX_5;DT(0,0,0)rec[-1]_rec[-8];TICK;RX_5;TICK;CX_5_9;TICK;CX_5_3_9_10;MARKX(0)3;TICK;CX_3_1_5_6_9_12;TICK;CX_1_0_9_8_6_7_12_11_4_5_2_3;MARKY(0)7;TICK;MX_12_9_4_2_6_1;S_5_11_7_0;S_DAG_3_8_10;DT(4,1,1)rec[-6];DT(3,1,1)rec[-5];DT(2,1,1)rec[-4]_rec[-7];DT(1,0,1)rec[-3];DT(2,2,1)rec[-2];DT(0,1,1)rec[-1];TICK;MPP_X8*X0*X3*X5;DT(3,0,2)rec[-1]_rec[-8]_rec[-14];TICK;MPP_X7*X10*X5*X3;DT(1,1,3)rec[-1]_rec[-9]_rec[-14];TICK;MPP_X11*X8*X5*X10;DT(4,0,4)rec[-1]_rec[-10]_rec[-14];TICK;MPP_Z8*Z0*Z3*Z5;DT(3,0,5)rec[-1]_rec[-14];TICK;MPP_Z7*Z10*Z5*Z3;DT(2,3,6)rec[-1]_rec[-14];TICK;MPP_Z11*Z8*Z5*Z10;DT(4,0,7)rec[-1]_rec[-14];TICK;MPP_Y0*Y8*Y11*Y10*Y7*Y5*Y3;OI(0)rec[-1]_rec[-8]_rec[-9]_rec[-12]_rec[-13]}{this Crumble link}
	to see that circuit with the error events from \cref{fig:2_fault_configuration}.}
	\label{fig:d3a6}
\end{SCfigure}

\begin{figure}[H]
	\centering
	\begin{tikzpicture}[
		stage/.style={
			draw=black,
			text centered,
			text depth=.25ex,
			text height=1.5ex
		},
		arrow/.style={-{Stealth[length=2mm]}},
	]
	% nodes
	\node[stage] (prepare) {prepare $\llbracket 7, 1, 3 \rrbracket$};
	\node[stage] (stabilise) [right=of prepare] {stabilise};
	\node[stage] (double_check) [right=of stabilise] {double check};
	\node[stage] (escape) [right=of double_check] {escape};
	\node[stage, yshift=3mm] (grow) [below=of $(stabilise)!0.5!(double_check)$] {grow distance by 2};
	% lines
	\draw[arrow] (prepare.east) -- (stabilise.west);
	\draw[arrow] (stabilise.east) -- (double_check.west);
	\draw[arrow] (double_check.east) -- (escape.west);
	\draw[arrow,dashed] (double_check.south) -- ([xshift=8mm]grow.north);
	\draw[arrow,dashed] ([xshift=-8mm]grow.north) -- (stabilise.south);
	\end{tikzpicture}
	\caption{Summary of magic state cultivation.
	Distance-5 cultivation traverses the dashed path once.}
	\label{fig:stage_flowchart}
\end{figure}

MSC-5,
also summarised in \cref{fig:stage_flowchart},
is the same as MSC-3 through to the distance-3 double check,
after which the $\llbracket 7, 1, 3 \rrbracket$ colour code
is grown into a $\llbracket 19, 1, 5\rrbracket$ colour code.
Three stabiliser measurement rounds are performed,
then $\overline{H_+}$ is again measured twice
using the distance-5 analogue of \cref{fig:d3a6};
we call this the \emph{distance-5 double check}.
Stage 3 is the same as for MSC-3---with
larger initial and final codes.
MSC-5 is intended to have fault distance 4 before its escape stage;
this increases to 5 if four instead of three stabiliser rounds
are performed before the distance-5 double check~\cite[\S 2.3]{Gidney2024} \cite[\S 4]{Hetenyi2026}.
Gidney et al.~\cite{Gidney2024} chose three instead of four rounds
because the spacetime cost of adding another round is not worth
the small improvement in logical error rate at practically relevant noise levels.

The Clifford analogue of T-state cultivation
is S-state cultivation,
whose circuit can be obtained by replacing all
T ($\T^\dag$) gates with S ($\S^\dag$).
As \cref{sec:introduction} mentions,
multiple papers~\cite{Gidney2024,Li2025,Tuloup2026,Chase2026,Fang2026a,Chen2026b}
report an unexplained difference in
the pre-escape logical error rate per kept shot
between T- and S-state cultivation,
which increases as the noise level decreases.